% \subsection*{2.4 极端暗光环境下融合偏振特征点的视觉自主导航}
\subsubsection*{2.4.2 暗光融合偏振特征点前端}

\textbf{(1) 偏振-可见光特征融合方法}

以上一节所述的暗光条件下的偏振特征提取方法为基础，在现有VINS-Fusion\cite{qin2019vinsfusion}前端框架的基础上，设计偏振-可见光特征融合方法，实现暗光融合偏振特征点前端。

在间接法视觉导航中，同一特征的持续跟踪帧数是影响系统精度和鲁棒性的关键因素。对同一三维空间点的连续观测构成了不同关键帧位姿之间的几何约束，观测的帧数越多，约束项越丰富，优化问题越不易受观测噪声和误匹配的干扰而发散，系统的导航精度和鲁棒性也随之提升。

然而，如图\ref{fig:polar_channels}所示，随着光照强度降低，可见光图像$S_0$信噪比急剧下降，特征跟踪极易中断。与此同时，偏振通道$P_0, P_1, P_2$在同等暗弱条件下仍能保持一定的结构信息。但由于各通道独立进行特征检测与光流跟踪，同一空间点可能在多个通道中被检测为独立的局部特征。若直接将这些特征不加区分地全部送入后端，不仅造成冗余，更关键的是，无法利用偏振通道在暗光下对同一空间点持续跟踪的能力来补偿$S_0$通道的特征丢失。因此，有必要设计偏振-可见光特征融合方法，将各通道独立跟踪的局部特征按照空间位置关联为统一的全局特征，从而在暗光条件下维持充足的可跟踪特征数量。

具体而言，对于同一帧偏振图像，偏振通道$P_0, P_1, P_2$与可见光图像$S_0$的相机位姿完全相同。物体边缘、角点等结构特征在不同通道的同一像素位置均有响应。这些不同通道检测到的同一空间位置的特征在物理上对应同一三维空间点，在逻辑上应合并为同一特征进行处理。

基于上述考虑，本文设计了跨通道全局特征池（Global Feature Pool），其核心思想是：为经空间位置关联后判定为同一空间点的各通道局部特征分配统一的全局标识符（Global ID），并以全局ID为单位向后端优化器提供观测。全局特征池的工作流程分为传播、关联和输出三个阶段，如图\ref{fig:global_feature_pool}所示。该流程图从左到右分为4个模块：(1) 每通道处理（Per-Channel Processing），展示4个输入通道$S_0$、$P_0(DoP)$、$P_1(\sin AoP)$、$P_2(\cos AoP)$各自独立进行特征点检测与光流跟踪；(2) 全局特征池（Global Feature Pool），包含传播（Propagation，检查上一帧全局特征关联的局部特征是否存活）、关联（Association，将新局部特征与已有全局特征按空间距离匹配）、输出（Output，每帧每全局特征仅输出一个通道的观测）三个子阶段；(3) 后端优化器（Bundle Estimator），接收全局特征观测并估计共享位姿。图中圆点表示局部特征，连线表示局部特征与全局特征的关联关系。

\begin{figure}[H]
\centering
\includegraphics[width=0.95\linewidth]{pic/1_polar_dark/7_global_feature_pool.png}
\caption{跨通道全局特征池工作流程示意图}
\label{fig:global_feature_pool}
\end{figure}

\textbf{(a) 传播阶段（Propagation）：}
对于上一帧中已存在的每个全局特征，遍历其关联的各通道局部标识符，检查这些局部特征在当前帧中是否仍被光流跟踪算法成功追踪。若某一通道的局部特征在当前帧中存活，则将该通道的当前观测（特征点像素坐标、归一化坐标等）重新绑定至该全局特征。这一机制使得全局特征的生命周期不再依赖于任一特定通道：只要至少有一个关联通道仍在跟踪，全局特征即可持续存在。

\textbf{(b) 关联阶段（Association）：}
对于当前帧中各通道新检测到或首次出现的局部特征，首先检查其是否可与已有的全局特征关联。具体地，计算该局部特征与当前帧中每个全局特征的最远关联点之间的像素距离，若该距离小于阈值$\tau_{\text{dist}}$（本文取$\tau_{\text{dist}}=4~\text{px}$，即特征点最小间距参数$\text{MIN\_DIST}=8$的一半），则将该局部特征合并至该全局特征。将$\tau_{\text{dist}}$设为$\text{MIN\_DIST}$的一半，是为了确保每个局部特征最多只与一个全局特征关联，避免出现同一局部特征被重复合并到多个全局特征的一对多问题。这一过程利用了多通道观测的空间一致性：由于各通道图像已通过偏振解码算法精确对齐，同一物理点在多个通道间的像素坐标偏差仅来源于滤波和量化误差，其大小通常不超过数个像素。若局部特征无法与任何已有全局特征关联，则为其创建新的全局特征。

\textbf{(c) 输出阶段（Output）：}
经传播和关联处理后，每个全局特征可能同时包含来自多个通道的观测。为避免向后端引入冗余信息，对于每个全局特征，每帧仅选择按照通道优先级第一个可用通道的观测结果（归一化坐标、像素坐标及速度）输出至后端优化器。后端优化器仅感知全局特征的存在，无需了解其来源于哪个通道或经历了几次跨通道关联，从而保持了对后端数据接口的完全兼容。

在上述输出阶段中，四个图像通道的通道优先级固定为：$S_0 \rightarrow P_0(DoP) \rightarrow$ $P_1(\sin AoP) \rightarrow P_2(\cos AoP)$，即优先使用可见光强度通道，其次为偏振度通道，最后为偏振角通道。该优先级的设计依据是：从前述滤波与特征提取实验可知，$S_0$与$DoP$通道在各类光照条件下的特征稳定性相对较高；而偏振角通道在计算过程中涉及$\arctan$运算，对$S_1,S_2$的噪声更为敏感，即使经过导向滤波后仍残留较多不确定性。因此，将偏振角通道作为高优先级通道的补充，仅在$S_0$和$DoP$均无法提供观测时启用。

需要指出的是，表\ref{tab:fp_metrics}数据显示，在暗光条件(0.9--4.8~lux)下，$P_2(\cos AoP)$通道的内点比例(0.650)优于$P_0(DoP)$的0.622（尽管内点数20.6仍低于$P_0$的37.2），表面上与上述优先级排序存在矛盾。对此，我们认为单通道特征质量并不等同于融合后的优先级——$P_0(DoP)$作为偏振度的标量表征，其空间分布与$S_0$强度信息的互补性更强，在全局特征池的关联阶段能够提供更多独立的空间约束；而$P_2$虽然内点比例较高，但其空间响应模式与$S_0/P_0$的重叠度更高，且内点绝对数量仅为$P_0$的约一半(20.6 vs 37.2)，作为高优先级输出时对全局特征数量的增量贡献有限。此外，优先级设计还考虑了计算效率因素：$P_0$的计算仅需Stokes矢量的简单运算，而$P_2$依赖$\arctan$中间结果，计算链路更长。尽管如此，当前固定优先级策略仍有改进空间，后续工作中我们将探索基于光照条件自适应的通道优先级策略，以进一步提升融合效果。

通过上述融合机制，偏振-可见光特征融合方法实现了以下优势：其一，暗光下可跟踪特征数量显著增加——即使在$S_0$通道因暗光噪声丢失大量特征后，偏振通道$P_0, P_1, P_2$仍能在相同空间位置检测到特征，通过全局特征池的关联机制补充为全局特征，使系统在极端暗光下仍可维持充足的全局特征数量；其二，后端优化问题的规模得到控制——合并后的特征数量不随通道数线性增长，避免了冗余观测对优化效率的影响；其三，多通道观测通过空间一致性约束在关联阶段即进行了隐式的交叉验证，一定程度上抑制了单通道中孤立误匹配进入后端。


\textbf{(2) 偏振融合特征分析}

为验证以上设计的偏振-可见光特征融合方法有效性，在实际采集的不同光照条件的数据上，运行上述偏振-可见光特征融合方法，进行量化对比分析。

具体来说，我们仍使用前述三种光照条件采集的数据，对四种可见光和偏振通道组合($S_0$/$P_0$/$P_1$/$P_2$、$S_0$/$P_0$、$S_0$/$P_1$/$P_2$、仅$S_0$)运行带有全局特征池的完整前端，在运行时逐帧统计以下指标：

\begin{itemize}
    \item \textbf{匹配融合特征点数}（Num Matched）：传播阶段结束后，当前帧中从上一帧成功跟踪的全局特征数量。该指标反映了系统可用的跨帧几何约束的丰富程度，直接决定后端优化的鲁棒性。
    \item \textbf{匹配率}（Match Ratio）：匹配融合特征点数与融合特征点总数的比值。该指标衡量了特征跟踪的连续性——匹配率越高，说明特征跨帧存活能力越强。
    \item \textbf{平均跟踪帧数}（Avg Track Length）：对匹配特征，遍历其绑定的各通道局部特征的$track\_cnt$，取最大值作为该全局特征的等效跟踪帧数，再对所有匹配特征求均值。该指标衡量特征被持续观测的时长，跟踪帧数越多，后端滑动窗口内的约束越充分。
    \item \textbf{有效观测数}（Effective Obs）：匹配特征点数与平均跟踪帧数的乘积。该综合指标同时考虑了"有多少特征"和"每个特征被跟踪多久"，是对后端可用约束总量的直观刻画。
\end{itemize}

上述指标统计结果与分析如下所示。

\textbf{(a) 全局指标分析}

在不同环境光照条件下，对各通道组合的整个轨迹的指标取均值，如图\ref{fig:plfp_bar}和表\ref{tab:polarfp_stats}所示。

\begin{table}[htbp]
\centering
\caption{不同光照条件下各通道组合的融合特征点统计指标均值}
\label{tab:polarfp_stats}
\begin{tabular}{lccccc}
\toprule
光照条件 & 通道组合 & 匹配数 & 匹配率 & 平均跟踪帧数 & 有效观测数 \\
\midrule
94--205 lux & $S_0$/$P_0$/$P_1$/$P_2$ & 530 & 0.844 & 22.92 & 12358 \\
94--205 lux & $S_0$/$P_0$ & 381 & 0.866 & 24.42 & 9496 \\
94--205 lux & $S_0$/$P_1$/$P_2$ & 466 & 0.837 & 22.68 & 10815 \\
94--205 lux & 仅$S_0$ & 263 & 0.892 & 26.56 & 7153 \\
\cmidrule{1-6}
2.6--18.6 lux & $S_0$/$P_0$/$P_1$/$P_2$ & 202 & 0.315 & 15.67 & 3113 \\
2.6--18.6 lux & $S_0$/$P_0$ & 141 & 0.304 & 18.02 & 2502 \\
2.6--18.6 lux & $S_0$/$P_1$/$P_2$ & 130 & 0.250 & 17.80 & 2267 \\
2.6--18.6 lux & 仅$S_0$ & 51 & 0.184 & 34.43 & 1634 \\
\cmidrule{1-6}
0.9--4.8 lux & $S_0$/$P_0$/$P_1$/$P_2$ & 97 & 0.134 & 6.83 & 678 \\
0.9--4.8 lux & $S_0$/$P_0$ & 63 & 0.130 & 8.87 & 563 \\
0.9--4.8 lux & $S_0$/$P_1$/$P_2$ & 54 & 0.094 & 6.92 & 406 \\
0.9--4.8 lux & 仅$S_0$ & 14 & 0.052 & 16.82 & 275 \\
\bottomrule
\end{tabular}
\end{table}

\begin{figure}[htbp]
\centering
\includegraphics[width=0.95\linewidth]{pic/1_polar_dark/8_plfp_bar.png}
\caption{不同光照条件下各通道组合的融合特征点统计指标}
\label{fig:plfp_bar}
\end{figure}

图\ref{fig:plfp_bar}为2×2的子图矩阵，4个子图分别为：左上为匹配特征点数（Num Matched），右上为匹配率（Match Ratio），左下为平均跟踪帧数（Avg Track Length），右下为有效观测数（Effective Obs）。每个子图中，横坐标为3种光照条件（94--205 lux、2.6--18.6 lux、0.9--4.8 lux），每组包含4个彩色柱状条，从左到右依次为：蓝色表示$S_0+DoP+AoP$（全部通道融合）、橙色表示$S_0+DoP$、绿色表示$S_0+AoP$、红色表示仅$S_0$。

根据表\ref{tab:polarfp_stats}和图\ref{fig:plfp_bar}，我们有如下观察：

首先，在平均跟踪特征点数量方面，在所有光强条件下，均为使用全部3个偏振通道增强时跟踪特征点数量最多。光照条件越差，使用偏振通道增强和不使用偏振通道增强的数量差距就越显著。在极弱光强条件下，偏振通道增强将前端跟踪的特征点数量提升了6倍，充分说明了在暗弱光环境下，偏振融合特征点是补充可见光图像视觉特征的有效补充。

其次，匹配率表示稳定跟踪特征点占全部特征点的比例。稳定跟踪的特征点比例越高，说明检测到的特征点质量越高、匹配效果越好。图表说明，在明亮光照条件下，所有方法的匹配率无明显区别，甚至仅$S_0$方法还略高。但一旦光照条件变差，偏振融合特征点的匹配率就显著高于仅$S_0$方法的匹配率。这说明，在暗弱光环境下，偏振融合特征点可以有效提升检测到的特征点质量，获得更好更多的匹配。

第三，有效观测数是视觉导航后端进行优化所使用的视觉信息多少的直观体现，它同时受跟踪的特征点数量和特征点的平均跟踪帧数的影响。虽然从平均跟踪帧数来看，各光照条件下都是仅$S_0$通道的平均跟踪帧数最大，但仅$S_0$通道的跟踪特征点数量过少，甚至可能无法满足视觉导航初始化的要求。根据表\ref{tab:polarfp_stats}和图\ref{fig:plfp_bar}，在各光照条件下，融合偏振通道的有效观测数都显著高于仅$S_0$，且所有偏振通道都使用的情况下有效观测数最多。这充分证明了偏振-可见光特征融合方法的优势。

\textbf{(b) 逐帧动态分析}

图\ref{fig:plfp_line_bright}至图\ref{fig:plfp_line_dark}均为1行3列的子图布局。从左到右3列依次为：匹配特征点数（Num Matched）、匹配率（Match Ratio）、有效观测数（Effective Obs）。横坐标为帧序号（Frame），纵坐标为对应指标值。每个子图中包含4条彩色曲线，分别对应4种通道组合：蓝色为$S_0+DoP+AoP$（全部通道融合）、橙色为$S_0+DoP$、绿色为$S_0+AoP$、红色为仅$S_0$。图\ref{fig:plfp_line_bright}为明亮光照条件，图\ref{fig:plfp_line_medium}为较暗光照条件，图\ref{fig:plfp_line_dark}为极暗光照条件。

\begin{figure}[H]
\centering
\includegraphics[width=0.95\linewidth]{pic/1_polar_dark/9_plfp_line_1.png}
\caption{明亮光照(94--205 lux)下融合特征点逐帧统计}
\label{fig:plfp_line_bright}
\end{figure}

\begin{figure}[H]
\centering
\includegraphics[width=0.95\linewidth]{pic/1_polar_dark/9_plfp_line_2.png}
\caption{较暗光照(2.6--18.6 lux)下融合特征点逐帧统计}
\label{fig:plfp_line_medium}
\end{figure}

\begin{figure}[H]
\centering
\includegraphics[width=0.95\linewidth]{pic/1_polar_dark/9_plfp_line_3.png}
\caption{极暗光照(0.9--4.8 lux)下融合特征点逐帧统计}
\label{fig:plfp_line_dark}
\end{figure}

从图\ref{fig:plfp_line_bright}可以看到，在明亮光照条件下，四种通道组合的匹配特征数均保持在较高水平（200--600个），匹配率稳定在0.8--1.0区间，有效观测数随场景运动呈现周期性波动（对应相机旋转/平移运动模式的变化）。在明亮场景下，偏振融合特征点可以有效提升匹配特征点数量和有效观测数量，不过对匹配率提升贡献不大。

从图\ref{fig:plfp_line_medium}中可以看到，随着光照强度降低，各通道的匹配特征数相比明亮光照条件显著下降，此时偏振融合特征点相比仅$S_0$的匹配特征数量明显增加，有效观测数量明显增加，说明偏振融合特征的有效性。


在图\ref{fig:plfp_line_dark}中可以看到，极暗光照条件下，偏振融合特征点与仅$S_0$的差异最为显著。仅$S_0$组合的匹配特征数在大部分帧中低于50个，部分帧中甚至降至20个以下。在这样的情况下，视觉导航系统可能已经难以维持足够的几何约束，面临跟踪丢失的风险。而使用全部偏振通道的偏振融合特征点的匹配特征数是仅$S_0$的数倍，有效观测数也明显高于仅$S_0$。这表明，在$S_0$通道几乎无法提供有效特征观测的极端暗光条件下，偏振通道通过全局特征池的跨通道融合机制，成功补偿了$S_0$通道的特征丢失，为后端优化提供了充足的几何约束。

综合上述量化对比分析可知，偏振-可见光特征融合方法实现了以下优势：

\begin{enumerate}
    \item \textbf{增加可跟踪特征数量}：在$S_0$通道因暗光噪声丢失大量特征后，偏振通道仍能在相同空间位置检测到特征，通过全局特征池的关联机制补充为全局特征，使系统在极端暗光下维持充足的全局特征数量；
    \item \textbf{提升后端约束丰富度}：有效观测数这一综合指标在三种光照条件下均表明，多通道组合相比单一$S_0$通道可为后端优化提供更充分的几何约束，且光照越暗，提升幅度越大；
    \item \textbf{控制优化规模}：通过空间位置关联将多通道局部特征合并为全局特征，避免了冗余观测，使后端优化问题规模不随通道数线性增长。
\end{enumerate}

\subsubsection*{2.4.3 导航实验验证}

为综合验证上一节所述的暗光融合偏振特征点方法，基于上节所设计的暗光融合偏振特征点前端与VINS-Fusion后端，集成暗光融合偏振特征点视觉惯性导航系统；采集不同光照条件下的室内暗光导航数据，在该数据集上进行测试，获取与现有基线方法的对比结果。

\textbf{(1) 真实数据采集}

为综合评估上述算法性能，并与现有基线方法进行对比，在室内采集包括单目偏振相机与IMU的连续视觉导航数据。

具体来说，我们使用的偏振相机型号为LUCID TRI050S-QC。该型号偏振相机为彩色偏振相机，实际使用时通过彩色转灰度，合并RGGB拜尔滤镜的彩色通道为灰度通道，作为灰度偏振相机使用；合并过程中通过插值恢复分辨率。相机原始数据分辨率为2448*2048。受限于相机采集的硬件条件，在相机采集数据时设置硬件降分辨率至1224*1024@7FPS。由于第一节所述偏振通道获取过程有超像素合并，实际获取的$S_0$/$P_0$/$P_1$/$P_2$通道分辨率均为612*512。实验时，为避免噪声影响，相机增益设置为0，曝光设置为自动曝光，但为保证数据实时性，曝光时间上限设置为140ms。

我们还使用了了6轴IMU。IMU型号为ADIS16465，数据获取频率为100Hz。偏振相机与IMU固定连接，集成到预先设计的硬件平台中(实际同时集成了双目偏振相机，但本实验仅使用左目数据)。相机内参与相机和IMU间的外参通过kalibr\cite{rehder2016kalibr}进行标定。通过搭载NVIDIA Jetson Xavier NX的小型计算机连接IMU与偏振相机，控制数据采集。在电气层面，偏振相机通过百兆以太网连接计算机，IMU通过RS232接口连接计算机。在数据采集时，相机和IMU同步获取数据，通过ROS1同步采集，保存为ROS1 bag格式。

数据采集的地点是北京航空航天大学新主楼D602实验室室内，采集日期为2026年4月16日。实验共采集不同光照条件下的8组数据，每组数据时长在99s到117s不等。通过开关实验房间内灯光和开关相邻房间不同灯组实现不同的光照条件。在每组数据采集前，使用光强计测量房间内各角度最小光强与最大光强，作为光照条件的量化标准。

数据采集场景如图 XXX 所示，各组数据具体情况如表 \ref{tab:dataset_stats} 所示。

\begin{table}[htbp]
    \centering
    \caption{不同光照条件下偏振视觉导航数据集统计信息}
    \label{tab:dataset_stats}
    \begin{tabular}{lccc}
        \toprule
        \multirow{2}{*}{\textbf{数据编号}} & \textbf{最小光强} & \textbf{最大光强} & \textbf{数据时长} \\
         & \textbf{(/lux)} & \textbf{(/lux)} & \textbf{(/s)} \\
        \midrule
        1 & 94 & 205 & 99 \\
        2 & 9.6 & 54.3 & 111 \\
        3 & 6.7 & 40.4 & 117 \\
        4 & 3.2 & 30.6 & 107 \\
        5 & 2.6 & 18.7 & 108 \\
        6 & 1.4 & 9.3 & 113 \\
        7 & 1.1 & 5.5 & 120 \\
        8 & 0.9 & 4.7 & 110 \\
        \bottomrule
    \end{tabular}
\end{table}

由于实验条件限制，数据采集环境中无法获取真值(Ground Truth)信息。为使各算法有可对比基准，通过在地面做标记的方式确保每次采集时起点与终点均相同，即每组轨迹的终点位置均为原点。导航算法以原点为算法起点位置，运行全轨迹后，取算法解算的终点导航位置，即为终端位置误差(Terminal Position Error, TPE)。后续实验中，均以TPE作为各算法评价指标。

需要说明的是，以TPE作为唯一评价指标存在以下局限性：首先，TPE仅反映轨迹终点的位置误差，完全无法反映轨迹中间过程的累积漂移情况。一条轨迹可能在运动过程中间段漂移数米，但最后恰好回到原点附近，此时TPE很小但实际导航已经失效。其次，地面标记的起点和终点一致性依赖人工操作，本身存在一定的对齐误差(约数厘米)，这构成了TPE测量的噪声基底。在后续工作中，我们计划使用激光测距仪、UWB定位系统或全站仪等设备获取关键位姿参考点，或使用成熟的离线SLAM系统(如ORB-SLAM3\cite{campos2021orbslam3}在完成回环检测优化后)生成伪真值轨迹，从而引入绝对轨迹误差(Absolute Trajectory Error, ATE)和相对位姿误差(Relative Pose Error, RPE)等更全面的评价指标，以弥补当前仅依赖TPE的方法论不足。


\textbf{(2) 基线方法对比}

我们选用VINS-Fusion\cite{qin2019vinsfusion}作为基线方法进行对比，使用TPE作为量化评价指标。基线方法仅使用偏振相机原始数据解算的$S_0$通道。PolarFP-VINS同时使用$S_0$/$P_0$/$P_1$/$P_2$四个通道。我们的后端实现和基线方法完全一致，排除了由于优化算法差异带来的影响。除此之外，PolarFP-VINS和基线方法的所有共同前后端参数均保持一致。基线方法和PolarFP-VINS都未使用回环检测模块，仅对比滑窗优化结果。

单目-惯性视觉里程计的空间尺度估计容易出现较大误差，而我们实验的环境很暗，视觉特征稀疏，初始化误差更大，加剧了这一问题。为避免尺度估计误差造成TPE对比失效，我们统一将所有轨迹都按照平面外接矩形最大边长方向放缩到统一尺度。我们实际测量了房间的最大长度(5.0m)，并采用该长度作为统一放缩基准。

需要指出的是，上述轨迹放缩操作本质上是一种后处理的尺度对齐，它掩盖了单目视觉里程计最核心的尺度漂移问题。由于PolarFP-VINS和VINS-Fusion在特征点数量、跟踪连续性和优化约束上存在差异，两者在滑窗优化中产生的尺度漂移程度可能不同，统一放缩后比较TPE的方式虽然能够反映轨迹形状的相对精度，但无法反映两种方法在绝对尺度估计上的差异。因此，本节的对比实验仅评估轨迹形状精度，不评估尺度估计精度。尺度估计精度的定量评估需要在后续工作中通过引入深度估计模块、RGB-D传感器或激光雷达等具有尺度感知能力的传感器予以解决。


本研究所提方法(PolarFP-VINS)与基线方法在上述采集的8组数据上的性能对比如表\ref{tab:performance_comparison}所示。其中，轨迹6的基线方法在第44秒后失去跟踪，随后轨迹完全发散。


\begin{table}[htbp]
\centering
\caption{不同算法在室内暗光环境下的性能对比}
\label{tab:performance_comparison}
% 调整列间距以适应长表头
\setlength{\tabcolsep}{4pt}
\begin{tabular}{l c c c c c}
\toprule
\textbf{数据} & \textbf{\makecell{最小\\光强\\/lux}} & \textbf{\makecell{最大\\光强\\/lux}} & \textbf{\makecell{数据时长\\/s}} & \textbf{\makecell{PolarFP-VINS\\TPE/m}} & \textbf{\makecell{VINS-Fusion\\TPE/m}} \\
\midrule
1 & 94 & 205 & 99 & 0.34 & \textbf{0.27} \\
2 & 9.6 & 54.3 & 111 & \textbf{1.50} & 2.22 \\
3 & 6.7 & 40.4 & 117 & 0.66 & \textbf{0.23} \\
4 & 3.2 & 30.6 & 107 & 0.74 & \textbf{0.44} \\
5 & 2.6 & 18.7 & 108 & \textbf{0.36} & 1.27 \\
6 & 1.4 & 9.3 & 113 & \textbf{0.09} & --(failed) \\
7 & 1.1 & 5.5 & 120 & \textbf{0.28} & 2.07 \\
8 & 0.9 & 4.7 & 110 & \textbf{1.19} & 2.70 \\
\bottomrule
\end{tabular}
\end{table}

从表\ref{tab:performance_comparison}中我们可以看到，在光照条件相对较好的4组数据(数据1-4)中，有三组数据(数据1/3/4)的结果，都是PolarFP-VINS终端位置误差大于基线方法。然而，在光照条件相对较差的4组数据(数据5-8)中，PolarFP-VINS的终端位置误差全部优于基线方法。这说明，PolarFP-VINS在最大光强20lux及以下的室内暗光条件下，导航精度优于基线方法。

然而，在明亮条件下(数据1/3/4)PolarFP-VINS精度反而低于基线方法，这一现象值得深入分析。结合前述特征提取实验数据和系统架构，我们认为可能原因包括：

首先，偏振通道在明亮条件下引入了冗余与干扰。当$S_0$通道已经包含丰富的结构信息时，偏振通道的结构信息与$S_0$高度重叠，全局特征池在关联阶段可能将同一物理点的多个通道响应误关联为不同的全局特征，引入空间关联误差。此外，明亮环境中物体表面的镜面反射产生的偏振模式更为复杂和不稳定，DoP和AoP的局部突变可能引入误特征，反而降低了整体跟踪精度。

其次，导向滤波在明亮条件下对偏振通道的特征过滤过于激进。从表\ref{tab:filter_metrics}数据可见，明亮光照(94--205 lux)下，导向滤波使$P_0$通道的内点数从无滤波的250.9降至152.7，虽然内点比例从0.868提升至0.910，但可用特征总量的大幅减少削弱了全局特征池的关联基础。在明亮条件下，偏振通道的噪声水平本身较低，过度滤波反而损失了有效特征。

第三，多通道融合链路增加了时序延迟累积。PolarFP-VINS需要经过四通道独立提取→全局特征池传播-关联-输出的完整链路，而VINS-Fusion仅需单一$S_0$通道的特征跟踪。在相机运动过程中，多通道间的处理时延差异可能导致空间关联时的像素坐标偏差累积，尤其在全局特征池的关联阶段，像素级对齐误差在运动状态下可能被放大。

上述分析表明，PolarFP-VINS的优势集中在暗光场景，而在明亮条件下存在性能回退的风险。后续工作中，考虑设计光照自适应的通道选择策略——在明亮条件下降低偏振通道权重或仅使用$S_0$通道，在暗光条件下自动提升偏振通道优先级，从而兼顾不同光照条件下的导航精度。


\begin{figure}[htbp]
\centering
\includegraphics[width=0.7\linewidth]{pic/1_polar_dark/10_baseline_comp_1.png}
\caption{最大光强30lux及以上数据导航轨迹与基线对比}
\label{fig:basline_comp_1}
\end{figure}

上述实验结果的轨迹图如图\ref{fig:basline_comp_1} 和图\ref{fig:basline_comp_2} 所示。每个子图上方标注了对应的数据编号（Trail [1]--Trail [8]）。图中红色曲线为本研究提出的方法（PolarFP-VINS），蓝色曲线为基线方法（VINS-Fusion）。每条曲线旁标注了该方法对应轨迹的放缩倍率。每条轨迹的终点处，圆形标记（\textbullet）表示算法解算的终点位置，三角形标记（$\blacktriangle$）表示真实终点位置（原点）。轨迹6的基线方法在第44秒后失去跟踪，随后轨迹完全发散，图中用"trunc 44s"标注，仅画出其前44秒的轨迹结果。

\begin{figure}[htbp]
\centering
\includegraphics[width=0.7\linewidth]{pic/1_polar_dark/10_baseline_comp_2.png}
\caption{最大光强20lux及以下数据导航轨迹与基线对比}
\label{fig:basline_comp_2}
\end{figure}

从图\ref{fig:basline_comp_1} 和图\ref{fig:basline_comp_2}中可以看到，在最大光强30lux及以上数据中，PolarFP-VINS总体终端误差精度不如基线方法。但在最大光强20lux及以下数据中，PolarFP-VINS可以检测更多的视觉特征点，在暗光环境下也能较好进行跟踪；而VINS-Fusion可跟踪的特征点数量以及有效观测数都相对较少，漂移较为严重；体现了设计的偏振融合特征策略的独特优势。

此外，还测试了两种当前最佳(State of the Art, SOTA)的直接法视觉导航方法：DSO\cite{engel2018dso}和D3VO\cite{yang2020d3vo}。DSO\cite{engel2018dso}是基于传统优化方法的直接法，D3VO\cite{yang2020d3vo}是使用深度学习方法估计深度等信息的直接法。两种方法在所有采集的数据上均无法完成整个轨迹解算。

事实上，直接法对环境光照有比较苛刻的要求，其光度不变假设要求在滑动窗口范围内，环境光度不能改变，或仅可按照预设的模型改变。在暗光环境下，视觉特征整体稀疏，且传感器噪声较大，直接法的光度不变假设无法被良好满足，因此无法有效进行视觉导航。这也说明，间接法在捕捉暗光环境下视觉特征点具有相对优势。

此外，本节实验数据和分析存在以下局限性：

首先，实验仅采集了8组数据、总时长约15分钟，所有数据均在同一地点(北京航空航天大学新主楼D602实验室)采集，缺乏场景多样性。不同房间的纹理丰富度、光照分布、空间尺度差异，以及不同运动模式(快速旋转、长距离平移、往复运动等)对视觉里程计的影响差异很大，当前数据量不足以全面覆盖这些变化。

其次，受实验设备硬件条件限制，本实验数据采集帧率为7FPS，低于主流视觉导航系统的典型工作帧率(20--30FPS)。虽然低帧率对间接法的影响相对可控(特征点关联的基线更长，几何约束更强)，但帧率降低意味着单位时间内可用观测减少，在快速运动时更容易发生特征丢失。因此，将本实验结论外推至更高帧率场景时需谨慎，PolarFP-VINS在20FPS以上帧率下的实时性能和精度表现有待进一步验证。

最后，7FPS下的前端处理时间分析(表\ref{tab:filter_metrics}和表\ref{tab:fp_metrics})仅报告了单通道单帧耗时，四通道并行运行的总前端耗时和全局特征池的传播-关联-输出流程耗时尚未统计，PolarFP-VINS相比VINS-Fusion的整体帧率下降幅度及其对实时性的影响仍需在后续实验中补充测量。

\textbf{(3) 消融实验-偏振通道分析}

\textbf{(4) 消融实验-特征点参数分析}

\textbf{(5) 对比实验-SuperPoint特征点}

\subsubsection*{2.4.4 小结}

面对极端暗弱光照条件下视觉特征稀疏导致视觉导航失效的问题，本研究利用偏振信息在暗光条件下仍能捕捉环境结构特征的优势，基于微阵列偏振相机获取偏振视觉信息；设计了基于偏振度和偏振角解算、$S_0$导向滤波和GFTT特征点的偏振特征提取方法，通过匹配率等指标的定量对比确认了设计有效性；设计了基于全局特征池的偏振融合特征点前端，通过匹配融合特征点数、有效观测数等定量指标对比确认了设计有效性；基于上述方法集成了PolarFP-VINS暗光偏振融合特征视觉导航方法，在真实暗光环境数据上与基线方法进行了对比，确认了所提方法在暗光条件下的有效性；通过偏振通道和特征点参数两个消融实验证明了多偏振通道的必要性。